\documentclass[11pt,a4paper]{article}
\usepackage{amsmath,amssymb,amsthm}
\usepackage[margin=2.5cm]{geometry}
\usepackage{hyperref}

\newtheorem{definition}{Definition}[section]
\newtheorem{theorem}[definition]{Theorem}
\newtheorem{proposition}[definition]{Proposition}
\newtheorem{lemma}[definition]{Lemma}
\newtheorem{corollary}[definition]{Corollary}
\newtheorem{conjecture}[definition]{Conjecture}
\newtheorem{remark}[definition]{Remark}

\title{Hyperseed Formalization:\\
  Bootleggers, Baptists, and the Coming AI Regulation}
\author{ProtomegaTron (automated formalization)}
\date{2026-09-18}

\begin{document}
\maketitle

\begin{abstract}
We formalize the intellectual structure of Ben Goertzel's Substack article
``Bootleggers, Baptists, and the Coming AI Regulation'' (2026-04-28)
within the Hyperseed ontology. The article applies Bruce Yandle's
regulatory capture theory to AI regulation: safety advocates (Baptists)
provide moral cover while Big Tech incumbents (Bootleggers) benefit from
compliance barriers. Key mappings: the Baptist/Bootlegger coalition as
section-fiber misalignment, regulatory capture as transition function
hijacking, compliance barriers as artificial fiber costs, concentration
as single-fiber fragility, the safety narrative as frozen section, and
pro-diversity regulation as fiber-promoting governance.
\end{abstract}

\tableofcontents

%% =================================================================
\section{Section-fiber misalignment}
\label{sec:misalignment}
%% =================================================================

\begin{theorem}[The Bootlegger/Baptist pattern --- claims 1--7, 21]
\label{thm:misalignment}
The Bootleggers and Baptists pattern: two groups support the same
policy for fundamentally different reasons. The Baptists (safety
advocates) support AI regulation out of genuine moral concern. The
Bootleggers (Big Tech incumbents) support the same regulations because
compliance costs create barriers to entry.

In Hyperseed terms, this is \emph{section-fiber misalignment}: the
policy has a visible section (moral justification, stated purpose)
that is misaligned with its structural fiber (actual effect on the
ecosystem):
\[
  s_{\text{stated}} = \text{``make AI safer''} \qquad \text{but} \qquad
  F_{\text{actual}} = \text{``concentrate AI development''}
\]

The section and fiber tell different stories. The section is what
gets debated in public (safety, responsibility, harm prevention).
The fiber is what actually happens when the policy is implemented
(barriers to entry, reduced competition, incumbent protection).

The structural insight: this misalignment isn't a conspiracy but a
\emph{structural property} of the regulatory process. The Bootleggers
don't need to conspire with the Baptists --- they just need to support
the same regulations. The section-fiber misalignment arises naturally
from the structure of regulatory politics, not from intentional
deception.

This is a general pattern: any time a policy has broad moral appeal
(compelling section) and concentrated economic beneficiaries (fiber
that favors specific actors), the Bootlegger/Baptist dynamic is
likely. The moral appeal provides political cover; the economic
benefit provides funding and lobbying power.
\end{theorem}

%% =================================================================
\section{Transition function hijacking}
\label{sec:capture}
%% =================================================================

\begin{proposition}[Regulatory capture --- claims 9--12, 22, 27]
\label{prop:capture}
The regulatory process is a transition function: it maps from the
current state (unregulated or lightly regulated) to a new state
(regulated). Regulatory capture means this transition function has
been hijacked --- it nominally serves the public interest but
actually maps to incumbent advantage.

Three mechanisms of hijacking:
\begin{enumerate}
\item \textbf{Compliance costs as fixed barriers.} Safety testing,
  documentation, auditing, and liability insurance create fixed costs
  that scale sublinearly with organization size. Large companies
  absorb them easily; small projects cannot afford them. The
  transition function maps ``new entrant'' to ``blocked.''

\item \textbf{Revolving door.} Regulators come from industry and
  return to industry. The regulations reflect industry's understanding
  of what's ``feasible'' --- which naturally favors existing approaches.
  The transition function is written by the entities it nominally
  regulates.

\item \textbf{Standard-setting capture.} When incumbents dominate
  standard-setting, the standards reflect their architecture, their
  tooling, their workflow. New fiber types (genuinely different AI
  approaches) are structurally excluded because they don't fit the
  standard's assumptions.
\end{enumerate}

In fiber bundle terms, standard-setting capture is a \emph{transition
function monopoly}: the incumbents control which transition functions
are ``valid'' (compliant), and only transition functions compatible
with their fiber structure qualify. New fiber types that would require
different transition functions are excluded not because they're
unsafe but because they're structurally incompatible with the
captured standards.
\end{proposition}

%% =================================================================
\section{Single-fiber fragility}
\label{sec:fragility}
%% =================================================================

\begin{theorem}[Concentration as safety risk --- claims 13--16, 24]
\label{thm:fragility}
The article's central irony: regulations meant to make AI safer may
make it more dangerous by reducing diversity and competition.
Concentrating AI development in a few large companies creates a
\emph{single-fiber system}: one dominant approach, one set of
assumptions, one failure mode.

In Hyperseed terms, concentration reduces the fiber bundle toward
a single-fiber system:
\[
  \text{Rich bundle: } E = \bigcup_{i=1}^{N} F_i \quad (N \text{ large})
  \qquad \xrightarrow{\text{regulation}} \qquad
  \text{Poor bundle: } E \approx F_1 \cup F_2 \cup F_3
\]

Single-fiber systems are fragile:
\begin{itemize}
\item \textbf{Correlated failures:} similar architectures fail in
  similar ways. A vulnerability in one system is likely present in
  all similar systems.
\item \textbf{Monoculture risk:} no diversity of approaches means
  no fallback if the dominant approach has fundamental limitations.
\item \textbf{Single point of capture:} a few organizations are
  easier to capture (by governments, by market forces, by internal
  misalignment) than a diverse ecosystem.
\end{itemize}

A rich fiber bundle (diverse ecosystem) is robust:
\begin{itemize}
\item \textbf{Uncorrelated failures:} different architectures fail
  differently. A vulnerability in one approach is unlikely to affect
  unrelated approaches.
\item \textbf{Multiple fallbacks:} if one approach has fundamental
  limitations, others can compensate.
\item \textbf{Distributed resilience:} many organizations with
  different governance structures resist capture better than a few
  similar ones.
\end{itemize}

The incumbents are not inherently safer: they have their own
misaligned incentives (profit maximization, market dominance,
shareholder value). Treating them as the ``responsible'' actors
is itself a section-fiber misalignment --- the section
(``responsible company'') doesn't match the fiber
(profit-driven corporation).
\end{theorem}

%% =================================================================
\section{Frozen section narrative}
\label{sec:narrative}
%% =================================================================

\begin{proposition}[The safety narrative --- claims 12, 25]
\label{prop:narrative}
The narrative ``AI is dangerous, only responsible large companies
should develop it'' serves both Baptist and Bootlegger interests
simultaneously. It's a compelling story that happens to have
concentration-promoting consequences.

In Hyperseed terms, this narrative is a \emph{frozen section}: a
particular framing of the situation, treated as if it were the
complete picture (the fiber itself), when in fact it captures only
one perspective:
\begin{itemize}
\item \textbf{What the section captures:} AI can cause harm; some
  development practices are careless; oversight is needed.
\item \textbf{What the section misses:} concentration itself is a
  risk; diversity promotes safety; incumbents have misaligned
  incentives; small players bring innovation.
\end{itemize}

This is the same structure as the Bible-as-frozen-section in the
Pope Leo article: a particular historical perspective (``big
companies are the responsible actors'') frozen as if it were an
eternal truth, when in fact it's a contingent framing that serves
particular interests.

The frozen section is difficult to challenge because it has genuine
moral content (the Baptist part). Criticizing the regulatory
framework sounds like criticizing safety itself --- which is exactly
the cover the Bootleggers need.
\end{proposition}

%% =================================================================
\section{Fiber-promoting governance}
\label{sec:alternative}
%% =================================================================

\begin{definition}[Pro-diversity regulation --- claims 17--20, 26]
\label{def:alternative}
The alternative: regulation that promotes diversity, competition, and
decentralization rather than concentration. Open standards,
interoperability requirements, anti-monopoly provisions, support for
small players.

In Hyperseed terms, this is \emph{fiber-promoting governance}:
governance that encourages the creation of new fibers rather than
protecting existing ones:
\begin{itemize}
\item \textbf{Open standards} = public transition functions. Anyone
  can implement them; no proprietary lock-in.
\item \textbf{Interoperability} = composable fibers. Different
  systems can interact through shared interfaces.
\item \textbf{Anti-monopoly} = fiber diversity maintenance. Actively
  prevent the bundle from collapsing to few fibers.
\item \textbf{Small-player support} = new fiber encouragement. Reduce
  barriers to creating new fibers, don't increase them.
\end{itemize}

This is the same structural principle as OpenBGI's fiber-compatible
governance and the Pope Leo article's fiber-level thinking: governance
that can accommodate diversity and evolution, not governance that
freezes the current state.

The key design criterion: regulation should be evaluated by its
\emph{fiber effect} (does it increase or decrease ecosystem
diversity?), not just by its \emph{section appeal} (does it sound
like it promotes safety?). Section-attractive, fiber-harmful
regulation is the Bootlegger/Baptist pattern in action.
\end{definition}

%% =================================================================
\section{Cross-references}
%% =================================================================

\begin{remark}[Connections to other formalizations]
\begin{itemize}
\item \textbf{Pope Leo \& Anthropic vs Teilhard} (2026-05-27): frozen
  sections. The safety narrative is a frozen section parallel to the
  Bible as frozen section --- both capture one perspective and treat
  it as complete.
\item \textbf{OpenBGI} (2026-05-07): fiber-promoting governance.
  OpenBGI's decentralized structure is the concrete alternative to
  the concentration that Bootlegger/Baptist regulation produces.
\item \textbf{The Anthropic Fable Farce} (2026-07-08): Anthropic as
  Bootlegger. Anthropic's RSP framework examined through the lens
  of regulatory capture.
\item \textbf{The Fable of Benevolent Throttling} (2026-07-12):
  defensive framing. The throttling regime embodies the
  Bootlegger/Baptist dynamic --- moral justification (safety) for
  structural concentration.
\item \textbf{How Decentralized AI Can Help Avert} (2026-06-19):
  decentralized safety. The structural argument that decentralization
  is safer than concentration.
\item \textbf{The Geopolitics of the Great AI Bet} (2026-08-08):
  geopolitical dimension. The Bootlegger/Baptist dynamic plays out
  internationally as well as domestically.
\end{itemize}
\end{remark}

\begin{conjecture}[Bootlegger/Baptist detection]
Conjecture: section-fiber misalignment in regulation can be detected
by checking whether the regulation's compliance costs scale
sublinearly with organization size. If they do, the regulation
disproportionately burdens small players --- a structural signature
of the Bootlegger/Baptist pattern regardless of the regulation's
stated intent.

More formally: for a regulation $R$, define the compliance cost
function $C_R(s)$ where $s$ is organization size. If
$C_R(s)/s$ is decreasing in $s$ (compliance cost per unit size
decreases with size), the regulation has concentration-promoting
fiber regardless of its stated section.

This test doesn't prove the regulation is bad --- some genuinely
safety-promoting regulations may have economies of scale. But it
flags regulations that require the Bootlegger/Baptist analysis:
is the concentration-promoting fiber an unavoidable consequence of
a genuine safety need, or is it the actual purpose disguised by a
safety section?
\end{conjecture}

\end{document}
